\section{Baire 纲定理及其推论}

\label{sec:1-2}

上一节解决了“极限该怎么说”，但这还不够。第 \ref{chap:00} 章里的第二个失效点是存在性：即便我们已经知道弱极限、子网和闭包该怎样表述，仍然需要某种结构来保证“坏集合不能处处出现”。这正是完备性和 Baire 范畴方法进入优化的地方。

Baire 方法之所以在泛函分析中占据中心地位，并不是因为它提供了某种旁门技巧，而是因为它把完备性转化成了一条结构结论：可数个过于稀薄的集合，不可能填满一个真正良好的空间。后文许多看似与“范畴”无关的结论，例如一致有界、开映射、闭图像以及典型参数上的稳定性，底层都在使用这一点。

\subsection{稀疏集、第一纲集与 Baire 空间}

\begin{definition}[Baire 范畴的基本概念]\label{def:baire-space}
设 $X$ 为拓扑空间，$A\subseteq X$。
\begin{enumerate}
  \item 若 $\operatorname{int}\overline{A}=\varnothing$，则称 $A$ 为\emph{无处稠密集}。
  \item 若 $A$ 可以写成可数个无处稠密集的并，则称 $A$ 为\emph{第一纲集}；否则称其为\emph{第二纲集}。
  \item 若 $R\subseteq X$ 的补集是第一纲集，则称 $R$ 为\emph{剩余集}。
  \item 若任意可数个稠密开集的交仍稠密，则称 $X$ 为\emph{Baire 空间}。
\end{enumerate}
\end{definition}

在 Baire 空间中，剩余集总是稠密的。因此“典型性质”常常被理解为“在一个剩余集上成立”。后文处理连续线性算子族、值函数正则性和参数稳定性时，这种语言会反复出现。

\subsection{完备度量空间中的 Baire 纲定理}

\begin{theorem}[完备度量空间上的 Baire 纲定理]\label{thm:baire-complete}
设 $(X,d)$ 为非空完备度量空间，$\{U_n\}_{n\ge 1}$ 为 $X$ 中一列稠密开集。则
\[
\bigcap_{n=1}^{\infty} U_n
\]
在 $X$ 中稠密。
\end{theorem}

\begin{proof}
任取非空开集 $V\subseteq X$。只需证明
\[
V\cap \bigcap_{n=1}^{\infty} U_n\neq \varnothing.
\]
因为 $U_1$ 稠密且开，故可在 $V\cap U_1$ 中选取闭球 $\overline{B}(x_1,r_1)$，满足
\[
\overline{B}(x_1,r_1)\subseteq V\cap U_1,\qquad r_1\le 2^{-1}.
\]
递归地，若已选得 $\overline{B}(x_n,r_n)\subseteq U_n\cap B(x_{n-1},r_{n-1})$ 且 $r_n\le 2^{-n}$，则由 $U_{n+1}$ 稠密开，可在 $U_{n+1}\cap B(x_n,r_n)$ 中选取闭球
\[
\overline{B}(x_{n+1},r_{n+1})\subseteq U_{n+1}\cap B(x_n,r_n),\qquad r_{n+1}\le 2^{-(n+1)}.
\]

于是得到一列递减的非空闭集
\[
F_n:=\overline{B}(x_n,r_n),\qquad F_{n+1}\subseteq F_n,\qquad \operatorname{diam}(F_n)\le 2r_n\to 0.
\]
任取 $y_n\in F_n$。对 $m>n$，有 $y_m,y_n\in F_n$，故
\[
d(y_m,y_n)\le \operatorname{diam}(F_n)\to 0.
\]
因此 $(y_n)$ 为 Cauchy 列，由完备性收敛到某个 $x\in X$。又因每个 $F_n$ 闭且对所有 $m\ge n$ 有 $y_m\in F_n$，故 $x\in F_n$ 对一切 $n$ 都成立。于是
\[
x\in F_1\subseteq V,\qquad x\in F_n\subseteq U_n\quad (n\ge 1),
\]
故 $x\in V\cap\bigcap_{n\ge 1}U_n$。
\end{proof}

\begin{example}[为什么完备性不能省略]\label{ex:baire-rationals}
在有理数空间 $(\mathbb Q,|\cdot|)$ 中，枚举 $\mathbb Q=\{q_n\}_{n\ge1}$ 并令 $U_n=\mathbb Q\setminus\{q_n\}$。每个 $U_n$ 都是稠密开集，但
\[
\bigcap_{n=1}^{\infty}U_n=\varnothing.
\]
把这个例子放在这里，是为了强调定理~\ref{thm:baire-complete} 的核心前提不是“可数交”，而是完备性本身。
\end{example}

\subsection{局部紧 Hausdorff 空间的版本}

\begin{lemma}[局部紧 Hausdorff 空间中的紧闭包开集]\label{lem:locally-compact-open}
设 $X$ 为局部紧 Hausdorff 空间，$O\subseteq X$ 为非空开集。则存在非空开集 $W$ 使得
\[
\overline{W}\subseteq O,\qquad \overline{W}\ \text{紧}.
\]
\end{lemma}

\begin{proof}
任取 $x\in O$。由局部紧性可取 $x$ 的开邻域 $V$ 使得 $\overline{V}$ 紧；由 Hausdorff 性，紧集是闭集。再以正则性选取开集 $W$ 满足
\[
x\in W\subseteq \overline{W}\subseteq V\cap O.
\]
于是 $\overline{W}$ 为紧集且包含于 $O$。
\end{proof}

\begin{theorem}[局部紧 Hausdorff 空间上的 Baire 纲定理]\label{thm:baire-locally-compact}
设 $X$ 为非空局部紧 Hausdorff 空间，$\{U_n\}_{n\ge1}$ 为 $X$ 中一列稠密开集。则 $\bigcap_{n\ge1}U_n$ 在 $X$ 中稠密。
\end{theorem}

\begin{proof}
任取非空开集 $V\subseteq X$。由引理~\ref{lem:locally-compact-open}，存在非空开集 $W_1$ 使
\[
\overline{W_1}\subseteq V\cap U_1,\qquad \overline{W_1}\text{ 紧}.
\]
递归地，若已构造出非空开集 $W_n$ 使 $\overline{W_n}$ 紧且 $\overline{W_n}\subseteq U_n\cap W_{n-1}$，则对非空开集 $W_n\cap U_{n+1}$ 再次应用引理~\ref{lem:locally-compact-open}，可得非空开集 $W_{n+1}$ 满足
\[
\overline{W_{n+1}}\subseteq W_n\cap U_{n+1},\qquad \overline{W_{n+1}}\text{ 紧}.
\]
于是 $\{\overline{W_n}\}_{n\ge1}$ 构成递减的非空紧集列。紧性保证有限交性质成立，因此
\[
\bigcap_{n=1}^{\infty}\overline{W_n}\neq\varnothing.
\]
取 $x$ 属于该交，则 $x\in \overline{W_1}\subseteq V$ 且对每个 $n$ 有 $x\in \overline{W_n}\subseteq U_n$。故
\[
x\in V\cap\bigcap_{n=1}^{\infty}U_n,
\]
从而交集稠密。
\end{proof}

\subsection{等价刻画与前瞻性推论}

\begin{proposition}[Baire 空间的等价刻画]\label{prop:baire-equivalent}
设 $X$ 为拓扑空间。下列命题等价：
\begin{enumerate}
  \item $X$ 是 Baire 空间；
  \item 任意第一纲集都没有内点；
  \item 任意剩余集都是稠密集；
  \item 任意非空开子集都是第二纲集。
\end{enumerate}
\end{proposition}

\begin{proof}
$(1)\Rightarrow(2)$：若 $M=\bigcup_{n\ge1}A_n$ 为第一纲集，其中每个 $A_n$ 无处稠密，则 $X\setminus \overline{A_n}$ 均为稠密开集。由 (1) 可知
\[
\bigcap_{n\ge1}(X\setminus\overline{A_n})
\]
稠密，于是其补集
\[
\bigcup_{n\ge1}\overline{A_n}
\]
没有内点。因为 $M\subseteq \bigcup_{n\ge1}\overline{A_n}$，故 $M$ 也没有内点。

$(2)\Rightarrow(3)$：若 $R$ 为剩余集，则 $X\setminus R$ 为第一纲集，由 (2) 知其内点为空，等价于 $R$ 稠密。

$(3)\Rightarrow(1)$：任意可数个稠密开集的交，其补为可数个闭无内点集之并，故是第一纲集。于是该交为剩余集，由 (3) 稠密。

$(2)\Rightarrow(4)$：若非空开集 $O$ 是第一纲集，则它作为自身的子集有非空内点，与 (2) 矛盾。

$(4)\Rightarrow(2)$：若第一纲集 $M$ 有非空内点 $O$，则 $O\subseteq M$，从而 $O$ 也是第一纲集，与 (4) 矛盾。
\end{proof}

\begin{corollary}[稠密 $G_\delta$ 集的典型性]\label{cor:dense-gdelta}
在 Baire 空间中，可数个稠密开集的交是稠密的 $G_\delta$ 集；因而任意剩余集都包含一个稠密 $G_\delta$ 子集。
\end{corollary}

\begin{proof}
前一结论来自定义；后一结论只需把剩余集写成第一纲集的补，再将对应的闭包补取交即可。
\end{proof}

\paragraph{对后文的前瞻}

\begin{remark}
开映射定理、闭图像定理与一致有界原理的共同底层机制，正是定理~\ref{thm:baire-complete} 与命题~\ref{prop:baire-equivalent}。本书在第 3 章讨论 Banach 空间工具箱时，将重新回到本节，不在这里提前展开线性空间版本的全部证明。
\end{remark}

\begin{example}[最小抽象例子：一致有界原理的拓扑影子]\label{ex:uniform-boundedness-shadow}
设 $\{T_i\}_{i\in I}$ 为 Banach 空间上的连续线性算子族。若对每个固定向量 $x$，集合 $\{\|T_i x\|:i\in I\}$ 有界，则“范数爆炸点集”可以写成可数个闭集的并。Baire 方法排除了这些闭集在整个空间中都太薄却又处处出现的可能，因此导出算子范数的一致有界性。把这个例子放在这里，是为了说明范畴方法不是抽象装饰，而是泛函分析主定理的发动机。
\end{example}

\begin{example}[有限维坍缩：欧氏空间中一切自动成为 Baire]\label{ex:finite-dimensional-baire}
任意欧氏空间都是完备度量空间，因此由定理~\ref{thm:baire-complete} 自动是 Baire 空间。把这个例子放在这里，是为了说明有限维线性代数中“满射线性映射自动开”“图像闭合行为良好”等现象，其实早已暗含在本节的完备性框架里，只是欧氏空间把这些前提隐藏得过于自然。
\end{example}

\begin{example}[应用触点：坏参数集很薄]\label{ex:thin-bad-parameters}
在优化模型的参数分析里，常见结论是“对典型参数，解映射满足某种正则性；例外参数至多构成第一纲集”。把这个例子放在这里，是为了给后文的对偶稳定性与值函数敏感性分析提供术语准备：所谓“典型”，往往并非指测度意义上的几乎处处，而是 Baire 范畴意义上的剩余集。
\end{example}

因此，本节回答的是第 \ref{chap:00} 章中的第二个问题：为什么存在性与稳定性论证需要某种“坏集合不能太大”的原理。再往前一步，优化和随机过程还要求空间既能承受完备性，又有足够好的可分性和 Borel 结构；这就把我们带到下一节的 Polish 空间。

\subsection*{有限维对照}

Boyd 书几乎不需要显式提到 Baire 范畴，因为有限维范数空间自动完备，且很多线性映射性质可由矩阵论直接验证。本节的价值在于指出：这些熟悉结论在无穷维里不再是线性代数常识，而是定理~\ref{thm:baire-complete} 与定理~\ref{thm:baire-locally-compact} 的推论。

\subsection*{习题（待补）}

\begin{exercise}[Baire 等价刻画占位]\label{exr:baire-equivalence-placeholder}
补一题：要求把命题~\ref{prop:baire-equivalent} 的四个条件补成完整循环证明，并解释“剩余”与“第一纲”之间的对偶关系。
\end{exercise}

\begin{exercise}[完备性桥接题占位]\label{exr:baire-complete-placeholder}
补一题：比较 $\mathbb R$ 与 $\mathbb Q$ 中稠密开集可数交的行为，说明为什么例~\ref{ex:baire-rationals} 恰好刻画了完备性的缺失。
\end{exercise}
